% 77-01(77쪽) — 01·02번 공통 그림 · 구간 (α, β) 에서의 y=f(x) 의 그래프(a, c 에서 위로 볼록한 꼭짓점 · b 에서 x축에 접하는 극소 · d 에서 아래로 볼록한 꼭짓점). 스캔 화질이 낮아 TikZ 로 다시 그림.
% 원문에 있는 것만: x축 라벨 α, a, b, c, d, β · 안내 점선 다섯 줄 · 라벨 y=f(x), O, x, y (눈금 값·좌표는 원문에 없다).
\begin{tikzpicture}[x=0.3cm, y=0.3cm, line width=0.6pt, font=\small]
  \draw[->, >=stealth, line width=0.7pt] (-4.9,0) -- (10.4,0) node[below=1pt] {$x$};
  \draw[->, >=stealth, line width=0.7pt] (0,-2.9) -- (0,4.8) node[left=1pt] {$y$};
  \draw[densely dashed, line width=0.4pt] (-3.65,0) -- (-3.65,-1.25);
  \draw[densely dashed, line width=0.4pt] (-1.15,1.3) -- (-1.15,0);
  \draw[densely dashed, line width=0.4pt] (3.3,2.95) -- (3.3,0);
  \draw[densely dashed, line width=0.4pt] (6.6,0) -- (6.6,-2.15);
  \draw[densely dashed, line width=0.4pt] (9.0,1.2) -- (9.0,0);
  \draw[line width=0.6pt] plot[smooth, tension=0.6] coordinates {%
    (-4.35,-2.75) (-4.0,-2.15) (-3.65,-1.25) (-3.3,-0.55) (-2.85,0) (-2.4,0.5)
    (-1.95,0.95) (-1.55,1.2) (-1.15,1.3) (-0.75,1.2) (-0.35,1.0) (0.1,0.7)
    (0.5,0.35) (0.85,0.12) (1.25,0) (1.65,0.12) (2.05,0.5) (2.5,1.3)
    (2.95,2.78) (3.3,2.95) (3.65,2.78) (4.05,2.1) (4.35,1.1) (4.6,0)
    (4.95,-0.9) (5.4,-1.65) (5.8,-1.95) (6.15,-2.08) (6.6,-2.15) (7.05,-2.08)
    (7.5,-1.75) (7.95,-1.05) (8.35,0) (8.7,0.7) (9.0,1.2) (9.3,1.9) (9.5,2.7)};
  \node[above, inner sep=1.5pt] at (-3.65,0) {$\alpha$};
  \node[below, inner sep=2pt] at (-1.15,0) {$a$};
  \node[below left=-2pt and -3pt] at (0,0) {$\pt{O}$};
  \node[below, inner sep=2pt] at (1.25,0) {$b$};
  \node[below, inner sep=2pt] at (3.3,0) {$c$};
  \node[above, inner sep=1.5pt] at (6.6,0) {$d$};
  \node[below, inner sep=2pt] at (9.0,0) {$\beta$};
  \node[anchor=west, inner sep=1pt] at (7.3,4.0) {$y=f(x)$};
\end{tikzpicture}
